\documentclass[11pt,a4paper]{article}
\usepackage[margin=24mm]{geometry}
\usepackage{amsmath,amssymb,amsthm}
\setlength{\emergencystretch}{2em}
\newtheorem{lemma}{Lemma}
\newtheorem{theorem}[lemma]{Theorem}
\newcommand{\gen}[1]{\langle#1\rangle}
\title{A crossing that exchanges the remaining two rows}
\author{Human Compression project: L044 continuation}
\date{}
\begin{document}\maketitle
\begin{lemma}[A pure alternating group from two rows]\label{alternating}
Let $t\ge3$ and let the distinguished points be
$b,a_0,a'_0,\ldots,a_{t-1},a'_{t-1}$. Suppose
\[
 q_i=(a_{i-1}\ a_i)(a'_{i-1}\ a'_i)\quad(1\le i<t),
 \qquad
 s=(b\ a_0)\prod_{i=1}^{t-1}(a_i\ a'_i)f,
\]
where $f$ is an involution supported outside these $2t+1$ points.
Then $\gen{s,q_1,\ldots,q_{t-1}}$ contains the pure natural
alternating group on all $2t+1$ distinguished points.
\end{lemma}
\begin{proof}
The permutation $c=(sq_1)^2$ is a pure five-cycle on
$T=\{b,a_0,a'_0,a_1,a'_1\}$: on $T$, the product $sq_1$
is a five-cycle, and off $T$ all factors have order at most
two and disjoint supports. Furthermore $q_1$ is a reflection
of that five-cycle and fixes $b$. Thus $D=\gen{c,q_1}$
is the natural dihedral group of order ten on $T$.

Choose its reflection $d$ fixing $a_1$. It moves $a'_1$
to a different point $z$ of $T\setminus\{a_1\}$. Since $d$
fixes $a_2,a'_2$, the product $q_2q_2^d$ is a pure three-cycle
on $\{a'_1,a'_2,z\}$: the common transposition $(a_1\ a_2)$
cancels. Conjugating this cycle by powers of $c$ gives
three-cycles on $\{a'_2,u,v\}$, where the pairs $\{u,v\}$
run through a connected five-cycle on $T$. Indeed every nonzero
step in the cyclic group of order five generates it.

For any connected graph on $T$, the three-cycles through a
fixed additional point and its edges generate the natural
alternating group on that enlarged set. To see this, take a
spanning tree and add its edges from a starting edge: each new
three-point alternating group overlaps the preceding support
in the additional point and one tree vertex, so the elementary
alternating-overlap lemma grows the entire alternating group.
We therefore obtain pure $\operatorname{Alt}(T\cup\{a'_2\})$.
Its conjugate by $q_2$ has a different six-point support with
five points in common. Their union is $T\cup\{a_2,a'_2\}$,
so overlap gives its pure $A_7$.

Successively conjugate by $q_3,\ldots,q_{t-1}$. At each step
exactly the next pair enters the union, while all but two points
of the old support remain in the intersection. These intersections
have size at least five. Repeated overlap gives the claimed
pure alternating group. The external factor $f$ never enters
this group: it was eliminated already in $(sq_1)^2$.
\end{proof}

\begin{theorem}[Exclusion of the row-swap crossing]
Retain the complete critical-head hypotheses, actual IP and all
intransitive ambient omissions. Suppose a later unique interior
cut has adjacent first shared label $g$ and crossing $j=g+1$.
Write its omission orbits as $C=C_g$ and $A$, with crossing
$(b\ a)$, $b\in C$. Suppose the reversed support equality has
\[
 |A|=2t,\quad t=r-g-1,\quad
 \operatorname{Fix}(P_{g-1})=A\cup\{b\},
 \qquad P_k=G_{\le k},
\]
and $Q=G_{>j}$ has two faithful natural $S_t$ orbits on $A$.
Let $v=\rho_g|_A$ exchange these two orbits equivariantly.
For $t\ge4$, the crossing $s=\rho_j$ fixes every point of
$A$ other than $a$. More precisely, it exchanges $a,b$ and
fixes $a'=a^v$ and every point of $A\setminus\{a,a'\}$.
The alternative in which it swaps the two remaining rows is
impossible, also when this alternative is specified for $t=3$.
\end{theorem}
\begin{proof}
Use $v$ to identify the two rows, and choose one point of
each pair in the first row. The pair action of $Q$ is $S_t$.
Its distinguished pair omissions are intransitive: if one
were transitive, adjoining $\rho_g$ would join the two rows
inside the corresponding ambient omission, making $A$ one
orbit there. Since the omitted mark is supported on $A$,
it could join no ambient orbits, contradicting full transitivity.
A selected edge per pair-image label is a forest: a cycle
would give an avoiding path for one selected omission crossing.
Its $t-1$ edges on $t$ vertices make it a tree. Deleting an
edge leaves exactly the two orbits of that omission, since the
remaining tree edges connect both sides. Each image preserves
all the other omission cuts; their common classes are singletons
and the endpoints of its own edge. Thus each image is precisely
that transposition. Nonadjacent labels have disjoint edges by
commutation, so no vertex has degree three and the tree is a path. The unique-crossing anchor is fixed by all
but the first $Q$ generator. Thus the actual $Q$ marking is
\[
 q_i=(a_{i-1}\ a_i)(a'_{i-1}\ a'_i),\quad
 a_0=a,\quad a'_0=a',\qquad1\le i<t.
\]
These facts concern the actual faithful row actions; no projected
intersection property is used.

The subgroup $Q_1=\gen{q_2,\ldots,q_{t-1}}$ has fixed set
exactly $C\cup\{a,a'\}$ when $t\ge3$. Since $s$ centralizes
$Q_1$, the unique crossing forces it to fix $a'$ and preserve
the remaining paired points. When $t\ge4$, the two nontrivial
$Q_1$-orbits are natural $S_{t-1}$ orbits of size at least three.
Their individual centralizers are trivial, and their equivariant
identification is unique. Thus the restriction of $s$ to their
union is either the identity or the full row swap.

Suppose it is the row swap, for any $t\ge3$. Then
\[
 s=(b\ a)\prod_{i=1}^{t-1}(a_i\ a'_i)f,
 \qquad \operatorname{supp}(f)\subseteq C\setminus\{b\}.
\]
Lemma~\ref{alternating} gives pure
$\operatorname{Alt}(A\cup\{b\})\le G_{\ge g+1}$.
The generator $\rho_g$ moves $b$, because $P_{g-1}$ fixes $b$
and $P_g$ is transitive on $C$. Set $c=b^{\rho_g}\in C\setminus\{b\}$.
It preserves $A$, so conjugation by $\rho_g$ and overlap give
\[
             \operatorname{Alt}(A\cup\{b,c\})\le G_{\ge g}.
\]
On the prefix side, $P_g$ has exterior image $\gen v$.
The crossing attaches $a$ and then $v$ attaches $a'$; both
generators preserve the other paired points. Hence
$C_{g+1}=C\cup\{a,a'\}$, and the complete prefix theorem
puts its pure alternating group in $P_{g+1}$.
Actual IP would therefore put pure
$\operatorname{Alt}(\{a,a',b,c\})$ in
\[
 P_{g+1}\cap G_{\ge g}=\gen{\rho_g,s}.
\]
This is impossible: a group generated by two involutions is
dihedral, with all odd-order elements in its cyclic rotation
subgroup, whereas $A_4$ has noncommuting three-cycles.
Only the identity restriction on the remaining rows survives
when $t\ge4$.
\end{proof}

\paragraph{Attribution and remaining scope.}
The source route is L044 Propositions 5.20--5.32. The acquired
pure-prefix and alternating-overlap proofs remain paid. The
five-cycle/reflection construction supplies an explicit new
pure alternating subgroup, and the final actual-IP obstruction
reuses the acquired noncommuting-three-cycle argument.
The identity restriction for $t\ge4$, the additional small
centralizer possibilities at $t=2,3$, other cuts, opposite heads
and terminal singleton remain unresolved here. No new computation,
formal theorem, whole-original or question credit is claimed.
\end{document}
